\section{Chevalley schemes}
\label{section:I.8}

\subsection{Related local rings}
\label{subsection:I.8.1}

For every local ring $A$, we denote by $\mathfrak m(A)$ the maximal ideal of
$A$.

\begin{lemma}[8.1.1]
\label{I.8.1.1}
Let $A$, $B$ be two local rings such that $A\subset B$; the following
conditions are equivalent:
\begin{enumerate}
  \item[(i)] $\mathfrak m(B)\cap A=\mathfrak m(A)$;
  \item[(ii)] $\mathfrak m(A)\subset\mathfrak m(B)$;
  \item[(iii)] $1$ does not belong to the ideal of $B$ generated by
    $\mathfrak m(A)$.
\end{enumerate}
\end{lemma}

It is evident that (i) implies (ii) and (ii) implies (iii); finally, if (iii)
holds, $\mathfrak m(B)\cap A$ contains $\mathfrak m(A)$ and does not contain
$1$, hence is equal to $\mathfrak m(A)$.

When the equivalent conditions of
(\hyperref[I.8.1.1]{8.1.1}) are fulfilled, we say that $B$ \emph{dominates}
$A$; it amounts to the same to say that the injection $A\to B$ is a
\emph{local} homomorphism. It is clear that, in the set of local subrings of a
ring $R$, the relation of domination is an order relation.

\begin{env}[8.1.2]
\label{I.8.1.2}
Consider now a \emph{field} $R$. For every subring $A$ of $R$, we denote by
$L(A)$ the set of local rings $A_\mathfrak p$, where $\mathfrak p$ runs
through the prime spectrum of $A$; they are identified with subrings of $R$
containing $A$. Since
$\mathfrak p=(\mathfrak p A_\mathfrak p)\cap A$, the map
$\mathfrak p\mapsto A_\mathfrak p$ from $\operatorname{Spec}(A)$ to $L(A)$
is bijective.
\end{env}

\begin{lemma}[8.1.3]
\label{I.8.1.3}
Let $R$ be a field and $A$ a subring of $R$. For a local subring $M$ of $R$ to
dominate a ring $A_\mathfrak p\in L(A)$, it is necessary and sufficient that
$A\subset M$; the local ring $A_\mathfrak p$ dominated by $M$ is then unique
and corresponds to $\mathfrak p=\mathfrak m(M)\cap A$.
\end{lemma}

Indeed, if $M$ dominates $A_\mathfrak p$, we have
$\mathfrak m(M)\cap A_\mathfrak p=\mathfrak p A_\mathfrak p$ by
(\hyperref[I.8.1.1]{8.1.1}), whence the uniqueness of $\mathfrak p$; on
the other hand, if $A\subset M$, then $\mathfrak m(M)\cap A=\mathfrak p$ is
prime in $A$, and since $A-\mathfrak p\subset M$, we have
$A_\mathfrak p\subset M$ and
$\mathfrak p A_\mathfrak p\subset\mathfrak m(M)$, hence $M$ dominates
$A_\mathfrak p$.

\oldpage[I]{165}
\begin{lemma}[8.1.4]
\label{I.8.1.4}
Let $R$ be a field, $M$, $N$ two local subrings of $R$, and $P$ the subring of
$R$ generated by $M\cup N$. The following conditions are equivalent:
\begin{enumerate}
  \item[(i)] There exists a prime ideal $\mathfrak p$ of $P$ such that
    $\mathfrak m(M)=\mathfrak p\cap M$,
    $\mathfrak m(N)=\mathfrak p\cap N$.
  \item[(ii)] The ideal $\mathfrak a$ generated in $P$ by
    $\mathfrak m(M)\cup\mathfrak m(N)$ is distinct from $P$.
  \item[(iii)] There exists a local subring $Q$ of $R$ dominating both $M$
    and $N$.
\end{enumerate}
\end{lemma}

It is clear that (i) implies (ii); conversely, if $\mathfrak a\ne P$, then
$\mathfrak a$ is contained in a maximal ideal $\mathfrak n$ of $P$, and since
$1\notin\mathfrak n$, $\mathfrak n\cap M$ contains $\mathfrak m(M)$ and is
distinct from $M$, hence $\mathfrak n\cap M=\mathfrak m(M)$, and similarly
$\mathfrak n\cap N=\mathfrak m(N)$. It is clear that if $Q$ dominates $M$ and
$N$, we have $P\subset Q$, and
\[
  \mathfrak m(M)=\mathfrak m(Q)\cap M
  =(\mathfrak m(Q)\cap P)\cap M,\qquad
  \mathfrak m(N)=(\mathfrak m(Q)\cap P)\cap N,
\]
hence (iii) implies (i); the converse is evident by taking
$Q=P_\mathfrak p$.

When the conditions of (\hyperref[I.8.1.4]{8.1.4}) are satisfied, we say,
following C.~Chevalley, that the local rings $M$ and $N$ are
\emph{related}.

\begin{proposition}[8.1.5]
\label{I.8.1.5}
Let $A$, $B$ be two subrings of a field $R$, and $C$ the subring of $R$
generated by $A\cup B$. The following conditions are equivalent:
\begin{enumerate}
  \item[(i)] For every local ring $Q$ containing $A$ and $B$, we have
    $A_\mathfrak p=B_\mathfrak q$, where
    $\mathfrak p=\mathfrak m(Q)\cap A$,
    $\mathfrak q=\mathfrak m(Q)\cap B$.
  \item[(ii)] For every prime ideal $\mathfrak r$ of $C$, we have
    $A_\mathfrak p=B_\mathfrak q$, where
    $\mathfrak p=\mathfrak r\cap A$, $\mathfrak q=\mathfrak r\cap B$.
  \item[(iii)] If $M\in L(A)$ and $N\in L(B)$ are related, they are
    identical.
  \item[(iv)] We have $L(A)\cap L(B)=L(C)$.
\end{enumerate}
\end{proposition}

Lemmas (\hyperref[I.8.1.3]{8.1.3}) and
(\hyperref[I.8.1.4]{8.1.4}) prove that (i) and (iii) are equivalent;
it is clear that (i) implies (ii), by applying it to $Q=C_\mathfrak r$;
conversely, (ii) implies (i), since if $Q$ contains $A\cup B$, it contains
$C$, and if $\mathfrak r=\mathfrak m(Q)\cap C$, we have
$\mathfrak p=\mathfrak r\cap A$ and $\mathfrak q=\mathfrak r\cap B$ by
(\hyperref[I.8.1.3]{8.1.3}). It is immediate that (iv) implies (i), since
if $Q$ contains $A\cup B$, it dominates a local ring
$C_\mathfrak r\in L(C)$ by (\hyperref[I.8.1.3]{8.1.3}); by hypothesis we
have $C_\mathfrak r\in L(A)\cap L(B)$, and
(\hyperref[I.8.1.1]{8.1.1}) and (\hyperref[I.8.1.3]{8.1.3}) prove that
$C_\mathfrak r=A_\mathfrak p=B_\mathfrak q$. Finally, let us prove that (iii)
implies (iv). Let $Q\in L(C)$; $Q$ dominates an $M\in L(A)$ and an
$N\in L(B)$ (\hyperref[I.8.1.3]{8.1.3}), hence $M$ and $N$, being
related, are identical by hypothesis. Since we then have $C\subset M$,
$M$ dominates a $Q'\in L(C)$ (\hyperref[I.8.1.3]{8.1.3}), hence $Q$
dominates $Q'$, which (\hyperref[I.8.1.3]{8.1.3}) necessarily implies
$Q=Q'=M$, hence $Q\in L(A)\cap L(B)$. Conversely, if
$Q\in L(A)\cap L(B)$, we have $C\subset Q$, hence
(\hyperref[I.8.1.3]{8.1.3}) $Q$ dominates a
$Q''\in L(C)\subset L(A)\cap L(B)$; $Q$ and $Q''$, being related, are
identical, hence $Q''=Q\in L(C)$, which completes the proof.

\subsection{Local rings of an integral scheme}
\label{subsection:I.8.2}

\begin{env}[8.2.1]
\label{I.8.2.1}
Let $X$ be an \emph{integral} prescheme and $R$ its field of rational
functions, identical to the local ring of the generic point $a$ of $X$; for
every $x\in X$, we know that $\mathscr{O}_x$ is canonically identified with a
subring of $R$ (\hyperref[I.7.1.5]{7.1.5}), and for every rational
function $f\in R$, the domain of definition $\delta(f)$ of $f$ is the open
set of $x\in X$ such that $f\in\mathscr{O}_x$. It follows
(\hyperref[I.7.2.6]{7.2.6}) that for every open subset $U\subset X$, we have
\begin{equation}
  \Gamma(U,\mathscr{O}_X)=\bigcap_{x\in U}\mathscr{O}_x.
  \tag{8.2.1.1}\label{I.8.2.1.1}
\end{equation}
\end{env}

\oldpage[I]{166}
\begin{proposition}[8.2.2]
\label{I.8.2.2}
Let $X$ be an integral prescheme and $R$ its field of rational functions.
For $X$ to be a scheme, it is necessary and sufficient that the relation
``$\mathscr{O}_x$ and $\mathscr{O}_y$ are related''
(\hyperref[I.8.1.4]{8.1.4}) between points $x$, $y$ of $X$ imply $x=y$.
\end{proposition}

Suppose this condition holds, and let us show that $X$ is separated. Let
$U$ and $V$ be two distinct affine open subsets of $X$, and $A$ and $B$ their
rings, identified with subrings of $R$; $U$ (resp. $V$) is therefore
identified (\hyperref[I.8.1.2]{8.1.2}) with $L(A)$ (resp. $L(B)$), and the
hypothesis implies (\hyperref[I.8.1.5]{8.1.5}) that, if $C$ is the subring
of $R$ generated by $A\cup B$, $W=U\cap V$ is identified with
$L(A)\cap L(B)=L(C)$. Moreover, we know
(\cite{I-1}, p.~5-03, prop.~4 \emph{bis}) that every subring $E$ of $R$ is
equal to the intersection of the local rings belonging to $L(E)$; $C$ is
therefore identified with the intersection of the rings $\mathscr{O}_z$ for
$z\in W$, in other words (\hyperref[I.8.2.1.1]{8.2.1.1}) with
$\Gamma(W,\mathscr{O}_X)$. Consider then the sub-prescheme induced by $X$
on $W$; to the identity homomorphism
$\varphi:C\to\Gamma(W,\mathscr{O}_X)$ there corresponds
(\hyperref[I.2.2.4]{2.2.4}) a morphism
$\Phi=(\psi,\theta):W\to\operatorname{Spec}(C)$; we will see that $\Phi$
is an \emph{isomorphism} of preschemes, from which it will follow that $W$ is
an \emph{affine} open subset. The identification of $W$ with
$L(C)=\operatorname{Spec}(C)$ shows that $\psi$ is \emph{bijective}. On the
other hand, for every $x\in W$, $\theta_x^\sharp$ is the injection
$C_\mathfrak r\to\mathscr{O}_x$, where
$\mathfrak r=\mathfrak m_x\cap C$, and by definition $C_\mathfrak r$ is
identified with $\mathscr{O}_x$, hence $\theta_x^\sharp$ is bijective. It
therefore remains to see that $\psi$ is a \emph{homeomorphism}, in other words
that for every closed subset $F\subset W$, $\psi(F)$ is closed in
$\operatorname{Spec}(C)$. Now, $F$ is the trace on $W$ of a closed subset of
$U$, of the form $V(\mathfrak a)$, where $\mathfrak a$ is an ideal of $A$;
let us show that $\psi(F)=V(\mathfrak a C)$, which will prove our assertion.
Indeed, the prime ideals of $C$ containing $\mathfrak a C$ are the prime
ideals of $C$ containing $\mathfrak a$, hence the ideals of the form
$\psi(x)=\mathfrak m_x\cap C$ where
$\mathfrak a\subset\mathfrak m_x$ and $x\in W$; since
$\mathfrak a\subset\mathfrak m_x$ is equivalent to
$x\in V(\mathfrak a)=W\cap F$ for $x\in U$, we indeed have
$\psi(F)=V(\mathfrak a C)$.

It follows that $X$ is separated, since $U\cap V$ is affine and its ring $C$
is generated by the union $A\cup B$ of the rings of $U$ and $V$
(\hyperref[I.5.5.6]{5.5.6}).

Conversely, suppose $X$ separated, and let $x$, $y$ be two points of $X$ such
that $\mathscr{O}_x$ and $\mathscr{O}_y$ are related. Let $U$ (resp.
$V$) be an affine open subset containing $x$ (resp. $y$), with ring $A$ (resp.
$B$); we then know that $U\cap V$ is affine and that its ring $C$ is generated
by $A\cup B$ (\hyperref[I.5.5.6]{5.5.6}). If
$\mathfrak p=\mathfrak m_x\cap A$, $\mathfrak q=\mathfrak m_y\cap B$, we have
$A_\mathfrak p=\mathscr{O}_x$, $B_\mathfrak q=\mathscr{O}_y$, and since
$A_\mathfrak p$ and $B_\mathfrak q$ are related, there exists a prime ideal
$\mathfrak r$ of $C$ such that $\mathfrak p=\mathfrak r\cap A$,
$\mathfrak q=\mathfrak r\cap B$ (\hyperref[I.8.1.4]{8.1.4}). But then
there exists a point $z\in U\cap V$ such that
$\mathfrak r=\mathfrak m_z\cap C$, since $U\cap V$ is affine, and we
evidently have $x=z$ and $y=z$, whence $x=y$.

\begin{corollary}[8.2.3]
\label{I.8.2.3}
Let $X$ be an integral scheme, and $x$, $y$ two points of $X$. For
$x\in\overline{\{y\}}$, it is necessary and sufficient that
$\mathscr{O}_x\subset\mathscr{O}_y$, in other words that every rational
function defined at $x$ be defined at $y$.
\end{corollary}

The condition is evidently necessary, since the domain of definition
$\delta(f)$ of a rational function $f\in R$ is open; let us show that it is
sufficient. If $\mathscr{O}_x\subset\mathscr{O}_y$, there exists a prime
ideal $\mathfrak p$ of $\mathscr{O}_x$ such that $\mathscr{O}_y$ dominates
$(\mathscr{O}_x)_\mathfrak p$ (\hyperref[I.8.1.3]{8.1.3}); now
(\hyperref[I.2.4.2]{2.4.2}) there exists $z\in X$ such that
$x\in\overline{\{z\}}$ and
$\mathscr{O}_z=(\mathscr{O}_x)_\mathfrak p$; since $\mathscr{O}_z$ and
$\mathscr{O}_y$ are related, we have $z=y$ by
(\hyperref[I.8.2.2]{8.2.2}), whence the corollary.

\begin{corollary}[8.2.4]
\label{I.8.2.4}
If $X$ is an integral scheme, the map $x\mapsto\mathscr{O}_x$ is
injective; in other words, if $x$, $y$ are two distinct points of $X$, there
exists a rational function defined at one of these points and not at the
other.
\end{corollary}

\oldpage[I]{167}
This follows from (\hyperref[I.8.2.3]{8.2.3}) and the axiom
$(T_0)$ (\hyperref[I.2.1.4]{2.1.4}).

\begin{corollary}[8.2.5]
\label{I.8.2.5}
Let $X$ be an integral scheme whose underlying space is Noetherian; when $f$
runs through the field $R$ of rational functions on $X$, the sets
$\delta(f)$ generate the topology of $X$.
\end{corollary}

Indeed, every closed subset of $X$ is then a finite union of irreducible
closed sets, that is, of sets of the form $\overline{\{y\}}$
(\hyperref[I.2.1.5]{2.1.5}). Now, if
$x\notin\overline{\{y\}}$, there exists a rational function $f$ defined at
$x$ and not at $y$ (\hyperref[I.8.2.3]{8.2.3}); in other words, we have
$x\in\delta(f)$ and $\delta(f)$ does not meet $\overline{\{y\}}$. The
complement of $\overline{\{y\}}$ is consequently a union of sets of the
form $\delta(f)$, and by virtue of the first remark, every open subset of $X$
is a union of finite intersections of open subsets of the form $\delta(f)$.

\begin{env}[8.2.6]
\label{I.8.2.6}
Corollary (\hyperref[I.8.2.5]{8.2.5}) shows that the topology of $X$ is
entirely characterized by the family of local rings
$(\mathscr{O}_x)_{x\in X}$ having $R$ as field of fractions. It amounts to
the same, moreover, to say that the closed subsets of $X$ are defined in the
following way: given a finite subset
$\{x_1,\ldots,x_n\}$ of $X$, we consider the set of $y\in X$ such that
$\mathscr{O}_y\subset\mathscr{O}_{x_i}$ for at least one index $i$, and these
sets (for all choices of $\{x_1,\ldots,x_n\}$) are the closed sets of $X$.
Moreover, once the topology of $X$ is known, the structural sheaf
$\mathscr{O}_X$ is also determined by the family of the
$\mathscr{O}_x$, since
$\Gamma(U,\mathscr{O}_X)=\bigcap_{x\in U}\mathscr{O}_x$
(\hyperref[I.8.2.1.1]{8.2.1.1}). The family
$(\mathscr{O}_x)_{x\in X}$ therefore completely determines the prescheme $X$
when $X$ is an integral scheme whose underlying space is Noetherian.
\end{env}

\begin{proposition}[8.2.7]
\label{I.8.2.7}
Let $X$, $Y$ be two integral schemes, $f:X\to Y$ a dominant morphism
(\hyperref[I.2.2.6]{2.2.6}), and $K$ (resp. $L$) the field of rational
functions on $X$ (resp. $Y$). Then $L$ is identified with a subfield of
$K$, and for every $x\in X$, $\mathscr{O}_{f(x)}$ is the unique local ring of
$Y$ dominated by $\mathscr{O}_x$.
\end{proposition}

Indeed, if $f=(\psi,\theta)$ and $a$ is the generic point of $X$, then
$\psi(a)$ is the generic point of $Y$
(\hyperref[0.2.1.5]{0, 2.1.5}); $\theta_a^\sharp$ is consequently a
monomorphism of the field $L=\mathscr{O}_{\psi(a)}$ into the field
$K=\mathscr{O}_a$. Since every non-empty affine open subset $U$ of $Y$ contains
$\psi(a)$, it follows from (\hyperref[I.2.2.4]{2.2.4}) that the
homomorphism
$\Gamma(U,\mathscr{O}_Y)\to\Gamma(\psi^{-1}(U),\mathscr{O}_X)$ corresponding
to $f$ is the restriction of $\theta_a^\sharp$ to
$\Gamma(U,\mathscr{O}_Y)$. Hence, for every $x\in X$,
$\theta_x^\sharp$ is the restriction of $\theta_a^\sharp$ to
$\mathscr{O}_{\psi(x)}$, and is consequently a monomorphism. We know moreover
that $\theta_x^\sharp$ is a local homomorphism; hence, if we identify $L$ with
a subfield of $K$ by $\theta_a^\sharp$, $\mathscr{O}_{\psi(x)}$ is dominated
by $\mathscr{O}_x$ (\hyperref[I.8.1.1]{8.1.1}); it is moreover the only
local ring of $Y$ dominated by $\mathscr{O}_x$, since two local rings of
$Y$ which are related are identical
(\hyperref[I.8.2.2]{8.2.2}).

\begin{proposition}[8.2.8]
\label{I.8.2.8}
Let $X$ be an irreducible prescheme and $f:X\to Y$ a local immersion
(resp. a local isomorphism); suppose moreover that the morphism $f$ is
separated. Then $f$ is an immersion (resp. an open immersion).
\end{proposition}

Let $f=(\psi,\theta)$; it suffices, in both cases, to prove that $\psi$
is a \emph{homeomorphism} of $X$ onto $\psi(X)$
(\hyperref[I.4.5.3]{4.5.3}). Replacing $f$ by $f_{\mathrm{red}}$
(\hyperref[I.5.1.6]{5.1.6} and
\hyperref[I.5.5.1]{5.5.1, (vi)}), we may suppose $X$ and $Y$ reduced. If
$Y'$ is the reduced closed sub-prescheme of $Y$ having
$\overline{\psi(X)}$ as underlying space, $f$ factors as
$X\xrightarrow{f'}Y'\xrightarrow{j}Y$, where $j$ is the canonical injection
(\hyperref[I.5.2.2]{5.2.2}). It follows from
(\hyperref[I.5.5.1]{5.5.1, (v)}) that $f'$ is again a separated morphism;
moreover, $f'$ is again

\oldpage[I]{168}
a local immersion (resp. a local isomorphism), since, the question being
local on $X$ and $Y$, we may restrict, in order to prove it, to the case where
$f$ is a closed immersion (resp. an open immersion), and our assertion then
follows immediately from (\hyperref[I.4.2.2]{4.2.2}).

We may therefore suppose that $f$ is a \emph{dominant} morphism, which
implies that $Y$ is irreducible as well
(\hyperref[0.2.1.5]{0, 2.1.5}), hence that $X$ and $Y$ are both
\emph{integral}. Moreover, the question being local on $Y$, we may suppose
that $Y$ is an affine scheme; since $f$ is separated, $X$ is a scheme
(\hyperref[I.5.5.1]{5.5.1, (ii)}), and we are finally under the
hypotheses of (\hyperref[I.8.2.7]{8.2.7}). Then, for every $x\in X$,
$\theta_x^\sharp$ is injective; but the hypothesis that $f$ is a local
immersion implies that $\theta_x^\sharp$ is surjective
(\hyperref[I.4.2.2]{4.2.2}), hence $\theta_x^\sharp$ is bijective;
in other words (with the identification of
(\hyperref[I.8.2.7]{8.2.7})), we have
$\mathscr{O}_{\psi(x)}=\mathscr{O}_x$. By
(\hyperref[I.8.2.4]{8.2.4}), this implies that $\psi$ is an
\emph{injective} map, which already proves the proposition when $f$ is a
local isomorphism (\hyperref[I.4.5.3]{4.5.3}). When we suppose only
that $f$ is a local immersion, for every $x\in X$ there exist an open
neighbourhood $U$ of $x$ in $X$ and an open neighbourhood $V$ of $\psi(x)$ in
$Y$ such that the restriction of $\psi$ to $U$ is a homeomorphism of $U$ onto
a \emph{closed} subset of $V$. Now, $U$ is dense in $X$, hence $\psi(U)$ is
dense in $Y$ and \emph{a fortiori} in $V$, which proves that
$\psi(U)=V$; since $\psi$ is injective, $\psi^{-1}(V)=U$, and this completes
the proof that $\psi$ is a homeomorphism of $X$ onto $\psi(X)$.

\subsection{Chevalley schemes}
\label{subsection:I.8.3}

\begin{env}[8.3.1]
\label{I.8.3.1}
Let $X$ be a \emph{Noetherian} integral scheme and $R$ its field of rational
functions; denote by $X'$ the set of local subrings
$\mathscr{O}_x\subset R$, where $x$ runs through $X$. The set $X'$ satisfies
the following three conditions:
\begin{enumerate}
  \item[(Sch. 1)] \emph{For every $M\in X'$, $R$ is the field of fractions
    of $M$.}
  \item[(Sch. 2)] \emph{There exists a finite set of Noetherian subrings
    $A_i$ of $R$ such that $X'=\bigcup_i L(A_i)$ and such that, for every pair
    of indices $i$, $j$, the subring $A_{ij}$ of $R$ generated by
    $A_i\cup A_j$ is an algebra of finite type over $A_i$.}
  \item[(Sch. 3)] \emph{Two elements $M$, $N$ of $X'$ which are related
    are identical.}
\end{enumerate}
\end{env}

Indeed, we have seen in (\hyperref[I.8.2.1]{8.2.1}) that (Sch. 1) is
satisfied, and (Sch. 3) follows from (\hyperref[I.8.2.2]{8.2.2}). To
prove (Sch. 2), it suffices to cover $X$ by finitely many affine open
subsets $U_i$ with Noetherian rings, and to take
$A_i=\Gamma(U_i,\mathscr{O}_X)$; the hypothesis that $X$ is a scheme implies
that $U_i\cap U_j$ is affine and that
$\Gamma(U_i\cap U_j,\mathscr{O}_X)=A_{ij}$
(\hyperref[I.5.5.6]{5.5.6}); moreover, since the space $U_i$ is
Noetherian, the immersion $U_i\cap U_j\to U_i$ is of finite type
(\hyperref[I.6.3.5]{6.3.5}), hence $A_{ij}$ is an $A_i$-algebra of finite
type (\hyperref[I.6.3.3]{6.3.3}).

\begin{env}[8.3.2]
\label{I.8.3.2}
The structures whose axioms are (Sch. 1), (Sch. 2) and (Sch. 3)
generalize the ``schemes'' in the sense of C.~Chevalley, who assumes moreover
that $R$ is an extension of finite type of a field $K$ and that the $A_i$ are
$K$-algebras of finite type (which makes part of (Sch. 2) superfluous)
\cite{I-1}. Conversely, if we have such a structure on a set $X'$, we can
associate with it an integral scheme $X$ by using the remarks of
(\hyperref[I.8.2.6]{8.2.6}): the underlying space of $X$ is equal to $X'$
endowed with the topology defined in (\hyperref[I.8.2.6]{8.2.6}), and with
the sheaf $\mathscr{O}_X$

\oldpage[I]{169}
such that $\Gamma(U,\mathscr{O}_X)=\bigcap_{x\in U}\mathscr{O}_x$ for every
open subset $U\subset X$, with an evident definition of the restriction
homomorphisms. We leave it to the reader to verify that we do indeed obtain
in this way an integral scheme, whose local rings are the elements of $X'$;
we will not use this result in what follows.
\end{env}
